\uuid{82pp}
\exo7id{4202}
\titre{exo7 4202}
\auteur{quercia}
\organisation{exo7}
\datecreate{2010-03-11}
\isIndication{false}
\isCorrection{true}
\chapitre{Equation différentielle}
\sousChapitre{Equations aux dérivées partielles}
\cpge{mpsi-mp2i.fonctions-deux-variables.b, pcsi.fonctions-deux-variables.b, ptsi.fonctions-deux-variables.b, mpsi-mp2i.fonctions-deux-variables.c, pcsi.fonctions-deux-variables.c, ptsi.fonctions-deux-variables.c, mpsi-mp2i.primitives-equations-differentielles.b, pcsi.techniques-calcul-primitives.b, ptsi.techniques-calcul-primitives.b, bcpst1.fonctions-deux-variables.c, bcpst1.derivees-primitives-integrales.b}
\module{Analyse}
\niveau{L2}
\difficulte{}

\contenu{
\question{
Déterminer les applications $f$ de classe $\mathcal{C}^1$ de ${\R^2}$ dans $\R$ vérifiant :
$\frac{\partial f}{\partial x} - \frac{\partial f}{\partial y} = a$ où $a$ est une constante réelle donnée.
On utilisera le changement de variable : $u = x+y$, $v = x-y$.
}
\reponse{
$f(x,y) = \frac{a(x-y)}2 + g(x+y)$.
}
}
